
\begin{definition}[Computational Complexity of Transformer Processing]
Let $L$ be the sequence length input to a transformer model. The computational complexity of the attention mechanism is $\mathcal{O}(L^2)$ and memory complexity is $\mathcal{O}(L^2)$.
\end{definition}

\begin{theorem}[Task Decomposition Efficiency]
Consider a task requiring total sequence length $L_{total}$ that can be decomposed into $n$ subtasks with sequence lengths $L_1, L_2, \ldots, L_n$ where $\sum_{i=1}^{n} L_i = L_{total}$.

The computational advantage of decomposition exists when:
\begin{equation}
\sum_{i=1}^{n} L_i^2 + C_{overhead} < L_{total}^2
\end{equation}

where $C_{overhead}$ represents coordination costs between agents.
\end{theorem}

\begin{proof}
The monolithic approach requires $\mathcal{O}(L_{total}^2)$ operations. The decomposed approach requires $\sum_{i=1}^{n} \mathcal{O}(L_i^2) = \mathcal{O}(\sum_{i=1}^{n} L_i^2)$ operations plus coordination overhead $C_{overhead}$.

Decomposition is beneficial when the total cost is less than monolithic processing:
$$\sum_{i=1}^{n} L_i^2 + C_{overhead} < L_{total}^2$$
\end{proof}

\begin{lemma}[Optimal Equal Division]
For equal task division where $L_i = \frac{L_{total}}{n}$ for all $i$, the computational savings factor is:
\begin{equation}
\frac{L_{total}^2}{\sum_{i=1}^{n} L_i^2} = \frac{L_{total}^2}{n \cdot \left(\frac{L_{total}}{n}\right)^2} = n
\end{equation}
\end{lemma}

\begin{proof}
Substituting $L_i = \frac{L_{total}}{n}$:
\begin{align}
\sum_{i=1}^{n} L_i^2 &= \sum_{i=1}^{n} \left(\frac{L_{total}}{n}\right)^2\\
&= n \cdot \frac{L_{total}^2}{n^2}\\
&= \frac{L_{total}^2}{n}
\end{align}

Therefore, the speedup factor is:
$$\frac{L_{total}^2}{\frac{L_{total}^2}{n}} = n$$
\end{proof}

\begin{theorem}[Memory Efficiency]
The memory savings for attention matrices under equal decomposition is:
\begin{equation}
\text{Memory Savings} = L_{total}^2 - \frac{L_{total}^2}{n} = L_{total}^2\left(1 - \frac{1}{n}\right)
\end{equation}
\end{theorem}

\begin{theorem}[Break-even Analysis]
Let $C_{overhead} = \alpha n + \beta n k + \gamma n m$ where:
\begin{itemize}
    \item $\alpha$ is the coordination cost per agent
    \item $\beta k$ is the communication cost per agent (k = message size)
    \item $\gamma m$ is the context switching cost per agent (m = model loading cost)
\end{itemize}

Task decomposition is beneficial when:
\begin{equation}
\frac{L_{total}^2}{n} + n(\alpha + \beta k + \gamma m) < L_{total}^2
\end{equation}

Rearranging:
\begin{equation}
L_{total}^2\left(1 - \frac{1}{n}\right) > n(\alpha + \beta k + \gamma m)
\end{equation}
\end{theorem}

\begin{corollary}[Optimal Number of Agents]
The optimal number of agents $n^*$ minimizes the total computational cost:
\begin{equation}
n^* = \arg\min_n \left\{\frac{L_{total}^2}{n} + n(\alpha + \beta k + \gamma m)\right\}
\end{equation}

Taking the derivative and setting to zero:
\begin{equation}
\frac{d}{dn}\left[\frac{L_{total}^2}{n} + n(\alpha + \beta k + \gamma m)\right] = -\frac{L_{total}^2}{n^2} + (\alpha + \beta k + \gamma m) = 0
\end{equation}

Solving for $n^*$:
\begin{equation}
n^* = \sqrt{\frac{L_{total}^2}{\alpha + \beta k + \gamma m}}
\end{equation}
\end{corollary}

\begin{theorem}[Quality-Efficiency Trade-off]
Let $A_{mono}$ be the accuracy of monolithic processing and $A_{decomp}$ be the accuracy of decomposed processing. The utility function is:
\begin{equation}
U = A_{decomp} \cdot \frac{L_{total}^2}{\sum_{i=1}^{n} L_i^2 + C_{overhead}} - \text{Cost}_{overhead}
\end{equation}

Decomposition is optimal when:
\begin{equation}
\frac{A_{decomp}}{A_{mono}} > \frac{\sum_{i=1}^{n} L_i^2 + C_{overhead}}{L_{total}^2}
\end{equation}
\end{theorem}

\begin{theorem}[Scalability Bound]
For tasks where sequence length grows as $L_{total} = \mathcal{O}(N^p)$ with problem size $N$, and optimal decomposition maintains $n = \mathcal{O}(N^q)$, the computational complexity scales as:
\begin{equation}
\mathcal{O}\left(\frac{N^{2p}}{N^q}\right) + \mathcal{O}(N^q) = \mathcal{O}(N^{2p-q}) + \mathcal{O}(N^q)
\end{equation}

For $q < 2p$, the first term dominates, giving complexity $\mathcal{O}(N^{2p-q})$.
\end{theorem}